%%%%%%%%%%%%%%%%%%%%%%%%%%%%%%%%%%%%%%%%%
% Medium Length Professional CV
% LaTeX Template
% Version 2.0 (8/5/13)
%
% This template has been downloaded from:
% http://www.LaTeXTemplates.com
%
% Original author:
% Trey Hunner (http://www.treyhunner.com/)
%
% Important note:
% This template requires the resume.cls file to be in the same directory as the
% .tex file. The resume.cls file provides the resume style used for structuring the
% document.
%
%%%%%%%%%%%%%%%%%%%%%%%%%%%%%%%%%%%%%%%%%

%----------------------------------------------------------------------------------------
%	PACKAGES AND OTHER DOCUMENT CONFIGURATIONS
%----------------------------------------------------------------------------------------

\documentclass{resume} % Use the custom resume.cls style
\usepackage[T1]{fontenc}
\usepackage[utf8]{inputenc}
\usepackage{array}
\usepackage{graphicx}
\usepackage[hidelinks]{hyperref}
\usepackage[left=0.75in,top=0.6in,right=0.75in,bottom=0.6in]{geometry} % Document margins
\newcommand{\tab}[1]{\hspace{.2667\textwidth}\rlap{#1}}
\newcommand{\itab}[1]{\hspace{0em}\rlap{#1}}
\name{Ziyan Xu} % Your name
\address{Globe Institute, University of Copenhagen \\ \O ster Voldgade 5-7 \\ 1350 Copenhagen K, Denmark} % Your address
%\address{123 Pleasant Lane \\ City, State 12345} % Your secondary addess (optional)
\address{ziyan.xu@sund.ku.dk \\ \href{http://ziyanxu.github.io}{ziyanxu.github.io} } % Your phone number and email

\begin{document}

\begin{rSection}{Education}

{\bf Peking University} \hfill {\em Beijing, China}
\\ Ph.D. in Astrophysics  \hfill {\em 2015 - 2021}
%\\ Kavli Institute for Astronomy and Astrophysics
%Minor in Linguistics \smallskip \\
%Member of Eta Kappa Nu \\
%Member of Upsilon Pi Epsilon \\

{\bf Peking University} \hfill {\em Beijing, China}
\\ B.S. in Astronomy, specializing in Astrophysics  \hfill {\em 2010 - 2015}
%\\ Department of Astronomy
\end{rSection}

\begin{rSection}{Research Experience}

\begin{rSubsection}{Marie Skłodowska-Curie Postdoctoral Fellow}{2025 -- present}{University of Copenhagen}{Supervisor: Prof. Anders Johansen}
\item \emph{Impact of dust evolution in new planet formation paradigm (PLANDEVOL)}
\end{rSubsection}

\begin{rSubsection}{Postdoctoral Researcher}{2024 -- 2025}{University of Copenhagen}{Supervisor: Prof. Anders Johansen}
\item \emph{{Pebble accretion and volatile transport in convective envelopes of protoplanets}}
\end{rSubsection}

\begin{rSubsection}{Postdoctoral Researcher}{2021 - 2024}{\'Ecole Normale Sup\'erieure de Lyon}{Supervisor: Prof. Guillaume Laibe}
\item \emph{{Local spherical collapsing box in Athena++}}\\
Collaborators: Prof. Guillaume Laibe, Dr. Elliot Lynch (CRAL - ENS Lyon)
\item \emph{{Dust dynamics and planetesimal formation in dust rings in turbulent protoplanetary disks}}\\
Collaborator: Prof. Xue-Ning Bai (Tsinghua University)
\item \emph{{Implementation of COALA high-order dust coagulation algorithm in Athena++}}\\
Collaborators: Dr. Maxime Lombart (National Taiwan Normal Univ.), Prof. Guillaume Laibe
\end{rSubsection}

\begin{rSubsection}{Ph.D. Researcher}{2015 - 2021}{Peking University}{Advisor: Prof. Gregory Herczeg}
%\item Advisor: Prof. Gregory Herczeg
\item \emph{{Probing protoplanetary disk wind with absorption line spectroscopy}}
\item \emph{{Dust dynamics in MRI turbulent protoplanetary disks}}\\
Collaborator: Prof. Xuening Bai (Tsinghua University)
\item \emph{{Atmospheric dynamics and circulation in warm Jupiters}}\\
Collaborator: Prof. Adam Showman (University of Arizona / Peking University)
\end{rSubsection}

\begin{rSubsection}{Student Intern}{2014 - 2015}{University of California, Santa Barbara}{Advisor: Prof. Ruth Murray-Clay}
%\item Advisor: Prof. Ruth Murray-Clay
\item \emph{{The effect of MRI turbulence on pebble accretion - data analysis}} 
%Analyzed simulation data and compared it with analytical theory.
\end{rSubsection}

\begin{rSubsection}{Student Researcher}{2014 - 2015}{Harvard-Smithsonian Center for Astrophysics}{Advisor: Prof. Xuening Bai}
%\item Advisor: Prof. Xuening Bai
\item \emph{{The effect of MRI turbulence on pebble accretion - simulation setup}} 
%Conducted pebble accretion simulations incorporating realistic disk turbulence using MHD simulations and evaluate the efficiency of the pebble accretion scenario of planetary core growth.
\end{rSubsection}


%\begin{rSubsection}{Peking University}{2012 - 2013}{Undergraduate Researcher}{}
%\item Advisor: Prof. Zuhui Fan
%\item Using n-body simulation to probe the formation and evolution of cosmic large-scale structures.
%\end{rSubsection}
%------------------------------------------------


\begin{rSubsection}{Visiting Observer (four observing nights)}{February 2019}{James Clerk Maxwell Telescope, Mauna Kea, Hawaii}{}
\item \emph{Four observing nights.}
\end{rSubsection}

\end{rSection}



\begin{rSection}{Fellowships, Grants and Awards} \itemsep -2pt

\begin{rSubsection}{Marie Skłodowska-Curie Postdoctoral Fellowship}{2025 -- 2027}{}{}
\item \emph{Impact of dust evolution in new planet formation paradigm (PLANDEVOL)}
\item Host: University of Copenhagen; Funding: EUR 247,553.28.
\end{rSubsection}

\item $\cdot$ \emph{Outstanding Graduates of Ordinary Higher Education Institutions} of Beijing (July 2021)
\item $\cdot$ \emph{Outstanding Graduates,} Peking University (July 2021)
\item $\cdot$ \emph{Specialty Scholarship (academic award),} Peking University (December 2019)
\item $\cdot$ \emph{Presidential Scholarship,} Peking University (December 2018)
\item $\cdot$ \emph{National Scholarship} of China (December 2017)
\item $\cdot$ \emph{Student Award for Outstanding Scientific Research,} Peking University (2016 \& 2017)
\item $\cdot$ \emph{May 4th Scholarship,} Peking University (December 2016)
\end{rSection}

\begin{rSectionLink}
{Publications}{https://ui.adsabs.harvard.edu/search/q=docs(library\%2Frj95jQ7nRomZnyOAC00MEg)\&sort=date\%20desc\%2C\%20bibcode\%20desc\&p_=0}{[ADS]}{\href{https://scholar.google.com/citations?user=Qf4qJ54AAAAJ}{[Google Scholar]}}\itemsep -3pt

\begin{list}{}{%
  \settowidth{\labelwidth}{2026}%
  \setlength{\labelsep}{0.6em}%
  \setlength{\leftmargin}{\dimexpr\labelwidth+\labelsep\relax}%
  \setlength{\itemsep}{0.65em}%
  \setlength{\parsep}{0pt}%
  \setlength{\topsep}{0.3em}%
}

\item[2026] Eriksson L. E. J., \textbf{Xu Z.}, Lim J., Yang C.-C., Huang P., \& Mac Low M.-M., \emph{Conditions for planetesimal formation via the streaming instability persist under turbulence driven by magnetorotational instability}, \href{https://ui.adsabs.harvard.edu/abs/2026A\%26A...714A..65E/abstract}{A\&A, 714, A65}

\item[2026] Wang H. S., Johansen A., \textbf{Xu Z.}, et al., \emph{Volatile depletion in rocky planets as a chemical fingerprint of hybrid accretion},
                                \href{https://ui.adsabs.harvard.edu/abs/arXiv:2608.24407}{Nature Astronomy}

\item[2025] Cui C. et al. (including \textbf{Xu Z.}), \emph{Dust Growth in ALMA Rings. II. Dusty Rossby Wave Instability}, \href{https://ui.adsabs.harvard.edu/abs/2025ApJ...986...86C/abstract}{ApJ, 986, 86}

\item[2024] \textbf{Xu Z.}, Lynch E. M., \& Laibe, G., \emph{Local spherical collapsing box in ATHENA++: Numerical implementation and benchmark tests}, \href{https://ui.adsabs.harvard.edu/abs/2024A\%26A...689A..96X/abstract}{A\&A, 689, A96}

\item[2023] Fang M. et al. (including \textbf{Xu Z.}), \emph{High-resolution [O I] line spectral mapping of TW Hya supportive of a magnetothermal wind}, \href{https://ui.adsabs.harvard.edu/abs/2023NatAs.tmp..121F/abstract}{Nature Astronomy, Volume 7, p. 905-912}

\item[2022] \textbf{Xu Z.} \& Bai X.-N., \emph{Turbulent Dust-trapping Rings as Efficient Sites for Planetesimal Formation}, \href{https://ui.adsabs.harvard.edu/abs/2022ApJ...937L...4X/abstract}{ApJL, 937, L4}

\item[2022] Espaillat C. C. et al. (including \textbf{Xu Z.}), \emph{The ODYSSEUS Survey. Motivation and First Results: Accretion, Ejection, and Disk Irradiation of CVSO 109}, \href{https://ui.adsabs.harvard.edu/abs/2022AJ....163..114E/abstract}{AJ, 163, 114}

\item[2022] \textbf{Xu Z.} \& Bai X.-N., \emph{Dust Settling and Clumping in MRI-turbulent Outer Protoplanetary Disks}, \href{https://ui.adsabs.harvard.edu/abs/2022ApJ...924....3X/abstract}{ApJ, 924, 3}

\item[2021] \textbf{Xu Z.}, Herczeg G. J., et al., \emph{Probing Protoplanetary Disk Winds with C II Absorption}, \href{https://ui.adsabs.harvard.edu/abs/2021ApJ...921..181X/abstract}{ApJ, 921, 181}

\item[2021] Lee Y.-H. et al. (including \textbf{Xu Z.}), \emph{The JCMT Transient Survey: Four-year Summary of Monitoring the Submillimeter Variability of Protostars}, \href{https://ui.adsabs.harvard.edu/abs/2021ApJ...920..119L/abstract}{ApJ, 920, 119}

\item[2017] \textbf{Xu Z.}, Bai X.-N., \& Murray-Clay R. A., \emph{Pebble Accretion in Turbulent Protoplanetary Disks}, \href{https://ui.adsabs.harvard.edu/abs/2017ApJ...847...52X/abstract}{ApJ, 847, 52}
\end{list}
\end{rSectionLink}

\begin{rSection}{Teaching and Supervision}
\begin{rSubsection}{Master's Internship Supervision}{2025}{University of Copenhagen / ENS de Lyon}{}
\item Supervised a Master's internship project.
\end{rSubsection}
\begin{rSubsection}{Teaching Assistant}{2019}{Peking University}{}
\item \emph{Modern Astronomy}, undergraduate course (Spring and Fall 2019).
\end{rSubsection}
\begin{rSubsection}{Teaching Assistant}{2017}{Peking University}{}
\item \emph{Stellar Structure and Evolution}, doctoral course (Fall 2017).
\end{rSubsection}
\end{rSection}

\begin{rSection}{Academic Service and Outreach}
\textbf{Academic service}
\begin{itemize}
\item Representative of the StarPlan division, Globe Postdoc Association (2026 -- present).
\item Local organising committee member, \emph{From Envelopes to Atmospheres}, Copenhagen (October 2026).
\item Organiser, weekly Planet Formation Meeting, University of Copenhagen (2025 -- present).
\item Referee, \emph{The Astrophysical Journal} and \emph{The Astrophysical Journal Letters} (2023 -- present).
\item Session moderator, ASIAA workshop (2021).
\item Volunteer, China-US Universities Astronomy Collaboration Summit (June 2016).
\end{itemize}

\textbf{Outreach and public engagement}
\begin{rSubsection}{StarPlan Outreach Coordinator, Copenhagen Culture Night}{October 2026}{University of Copenhagen, Faculty of Health and Medical Sciences}{}
\item Initiated StarPlan's participation, assembled the initial team, submitted the proposal, and served as the main contact with event organisers.
\item Coordinated collaborative programme development and led the Universe Sandbox activity at the team's public outreach booth.
\end{rSubsection}
\begin{itemize}
\item Speaker, Tycho Science Day, University of Copenhagen (2024).
\item Speaker and panel participant, International Day of Women and Girls in Science, Globe Diversity Programme (2024).
\end{itemize}
\end{rSection}

\begin{rSection}{Selected Conference Presentations \& Seminars} \itemsep -2pt
\begin{rSubsection}{From Envelopes to Atmospheres, Copenhagen, Denmark}{October 2026}{}{}
\item Contributed talk, \emph{"Pebble Accretion in Convective Envelopes: The Role of Dust Size and Aerodynamic Drag"}
\end{rSubsection}

\begin{rSubsection}{Discs on the Exe: From protostellar discs to debris discs and planets, Exeter, UK}{July 2026}{}{}
\item Contributed talk, \emph{"Runaway planetesimal formation in turbulent dust rings driven by a dust growth--concentration feedback loop"}
\end{rSubsection}

\begin{rSubsection}{SPiCE 2: Simulating Physics in Celestial Ecosystems, Lyon, France}{March 2026}{}{}
\item Invited talk, \emph{"Dust Evolution in Multiple Planet Formation Processes"}
\end{rSubsection}

\begin{rSubsection}{Flatiron Workshop on Hydrodynamic and Dusty Turbulence in Protoplanetary Disks, New York, NY, USA}{February 2026}{}{}
\item Contributed talk, \emph{"Feedback-regulated Dust Diffusion and Runaway Planetesimal Formation in Turbulent Disk Rings"}
\end{rSubsection}

\begin{rSubsection}{From Galaxies to Planets: An Elemental Journey, Ringberg Castle, Germany}{February 2026}{}{}
\item Invited talk, \emph{"Pebble Accretion and Volatile Transport in Convective Envelopes of Protoplanets"}
\end{rSubsection}

\begin{rSubsection}{International Conference on Exoplanets and Planet Formation (EPF 2025), Shanghai, China}{December 2025}{}{}
\item Poster, \emph{"Pebble Accretion and Volatile Transport in Convective Envelopes of Protoplanets"}
\item Poster, \emph{"Feedback-Regulated Dust Diffusion and Runaway Planetesimal Formation in Protoplanetary Disk Rings"}
\end{rSubsection}

\begin{rSubsection}{University of Arizona, AZ}{April 2025}{}{}
\item Steward/LPL Origins Seminar Talk, \emph{From Pebbles to Planets: Dust Dynamics in Planet Formation Processes} 
\end{rSubsection}
\begin{rSubsection}{Planet Formation \& Migration near Inner Edge of Disks, Santa Barbara}{April 2025}{}{}
\item Joint review talk, \emph{Modeling the Density and Thermal Structure near the Inner Rim} 
\end{rSubsection}
\begin{rSubsection}{Simulating Physics in Celestial Ecosystems, Sendai, Japan}{April 2024}{}{}
\item Contributed talk, \emph{Microphysics in Star \& Protoplanet Formation with Local Simulations} 
\end{rSubsection}
\begin{rSubsection}{Athena++ Workshop, New York, NY}{May 2023}{}{}
\item Contributed talk, \emph{Dust Coagulation in Athena++} 
\item Poster, \emph{Local Collapsing Boxes} 
\end{rSubsection}
\begin{rSubsection}{Protostars and Planets VII, Kyoto, Japan}{April 2023}{}{}
\item Poster, \emph{Turbulent Dust-trapping Rings as Efficient Sites for Planetesimal Formation} 
\end{rSubsection}
\begin{rSubsection}{MIAPP Workshop, Garching, Germany}{June 2022}{}{}
\item Contributed talk, \emph{Dust Dynamics and Planetesimal Formation in Ring-like Disk Substructures} 
\end{rSubsection}
\begin{rSubsection}{Lorentz Center Workshop, (Virtual)}{September 2021}{}{}
\item Contributed talk, \emph{Dust Dynamics and Planetesimal Formation in Turbulent Protoplanetary Disks and Dust Rings} 
\end{rSubsection}
\begin{rSubsection}{National Conference of Planetary Science, Suzhou, China}{June 2021}{}{}
\item Contributed talk, \emph{Dust Dynamics and Implications for Planetesimal Formation in Turbulent Protoplanetary Disks} 
\end{rSubsection}
\begin{rSubsection}{Five Years after HL Tau: a New Era in Planet Formation, (Virtual)}{December 2020}{}{}
\item Poster, \emph{Dust Dynamics and Planetesimal Formation in MRI Turbulent Protoplanetary Disks} 
\end{rSubsection}
\begin{rSubsection}{5th Workshop of the Network of Ultraviolet Astronomy, (Virtual)}{October 2020}{}{}
\item Contributed talk, \emph{Probing Protoplanetary Disk Winds with FUV Absorption Lines} 
\end{rSubsection}
\begin{rSubsection}{Planet Formation Workshop, Tokyo, Japan}{November 2019}{}{}
\item Contributed talk, \emph{Dust Settling and Clumping in MRI Turbulent Protoplanetary Disks} 
\end{rSubsection}
\begin{rSubsection}{Great Barriers in Planet Formation, Palm Cove, Australia}{July 2019}{}{}
\item Contributed talk, \emph{Probing Protoplanetary Disk Wind with Absorption Line Spectroscopy} 
\item Poster, \emph{Dust Settling and Clumping in MRI Turbulent Protoplanetary Disks} 
\end{rSubsection}
\begin{rSubsection}{Formation and Evolution of Solar System and Exoplanetary Systems, Urumqi, China}{July 2019}{}{}
\item Poster, \emph{Dust Settling and Clumping in MRI Turbulent Protoplanetary Disks}
\end{rSubsection}
\begin{rSubsection}{East Asian Observatory, Hilo, Hawaii}{March 2019}{}{}
\item EAO Seminar Talk, \emph{Absorption Lines as a Possible Probe of Protoplanetary Disk Wind} 
\end{rSubsection}
\begin{rSubsection}{Astrochemistry: Past, Present and Future, Caltech, CA}{July 2018}{}{}
\item Poster, \emph{C II Absorption Lines as a Possible Probe of Disk Photoevaporative Wind} 
\end{rSubsection}
\begin{rSubsection}{University of New South Wales, Sydney, Australia}{March 2018}{}{}
\item Astro Seminar Talk, \emph{Physical Processes in Protoplanetary Disks: Pebble Accretion and Disk Photoevaporation} 
\end{rSubsection}
\begin{rSubsection}{Exoplanets and Planet Formation, Shanghai, China}{December 2017}{}{}
\item Poster, \emph{Pebble Accretion in Turbulent Protoplanetary Disks} 
\end{rSubsection}
\begin{rSubsection}{OWL Summer Program, University of California, Santa Cruz, CA}{July 2017}{}{}
\item Contributed talk, \emph{Pebble Accretion in Turbulent Protoplanetary Disks} 
\end{rSubsection}
\begin{rSubsection}{Stanford University, CA}{July 2017}{}{}
\item KIPAC Tea Talk, \emph{Pebble Accretion in Turbulent Protoplanetary Disks} 
\end{rSubsection}
%\begin{rSubsection}{Earlier Conferences Attending}{}{}{}
%\item Chinese Astronomical Society Annual Meeting, Peking University (September 2015)
%\item The Disk in Relation to The Formation of Planets And Their Protoatmospheres Workshop, International Space Science Institute - Beijing (August 2014)
%\item Spring School on Protoplanetary Disks and Planet Formation, Tsinghua University (May 2014)
%\end{rSubsection}


\end{rSection}

\begin{rSection}{Technical Skills and Languages}

\begin{tabular}{ @{} >{\bfseries}l @{\hspace{6ex}} l }
Astrophysical & Athena/Athena++, RADMC-3D, MITgcm, VisIt, DS9, IRAF \\
Programming &  Python, C/C++, Fortran, MATLAB, shell, LaTeX\\
Language Skills & Mandarin (native), English (fluent), French (basic) \\
\end{tabular}

\end{rSection}

\end{document}
